%==============================================================================
% Chapter 5 — System Implementation & Testing
% سبرينت الفصل الخامس: 15/11 → 14/12 (حزمة جانت ٢.٣ — 30 يومًا — معلم M10)
% سير العمل: كل قسم يُكتب من مالكه بعد اكتمال مكوّنه في مستودع الكود،
% عبر فرع خاص ثم Pull Request يراجعه القائد (الفرع main محمي بقواعد
% إلزامية: لا دفع مباشر، وفحص بناء PDF إلزامي قبل الدمج).
% المراجعة الختامية والاعتماد: 12/12 → 14/12 (مهمة جانت ٢.٣.٩).
%==============================================================================
\chapter{تنفيذ واختبار النظام}
\label{ch:implementation}
\markboth{الفصل الخامس: تنفيذ واختبار النظام}{الفصل الخامس: تنفيذ واختبار النظام}

%------------------------------------------------------------------------------
% 5.1 بيئة التطوير والأدوات
% المالك: وليد العنسي — قابل للكتابة فور انطلاق السبرينت
% الموعد (السبرينت): 15/11 → 18/11
% المصادر المعتمدة:
%   (1) الفصل الثاني: الأدوات المختارة ومبررات اختيارها
%   (2) القسم 4.1: المكونات المعمارية وأصداراتها
%   (3) مستودع الكود Bayan: docker-compose.yml وMakefile وقرار البنية
%       التحتية ADR-05 (محليًا أولًا + GPU بالساعة للتقييم)
%   (4) مستودع openeducat-yemen-university: بيئة Odoo 17 + OpenEduCat
% المخرجات المتوقعة: جدول الأدوات بالأصدارات (اللغات/الأطر/قواعد البيانات/
%   النماذج اللغوية/الخدمات) مع مرجعية كل اختيار لما سبق، وجدول عتاد
%   التطوير والتشغيل، ووصف بيئة الحاويات الموحدة وبرنامج تشغيلها
%------------------------------------------------------------------------------
\section{بيئة التطوير والأدوات}
\label{sec:dev-env}
~

%------------------------------------------------------------------------------
% 5.2 تنفيذ النظام — يُكتب تدريجيًا مع اكتمال كل مكوّن في مستودع الكود
% التنظيم الداخلي المقترح (أقسام فرعية تُنشأ عند الكتابة):
%   5.2.1 تنفيذ منظومة الوكلاء — المالك: أحمد المتوكل — 15/11 → 08/12
%         (د١ وكيل SQL: مهمة ٢.٣.١؛ د٢ وكيل البحث: مهمة ٢.٣.٢؛
%          المشرف والمحلل والمُعِدّ: مهمة ٢.٣.٤ الموازية)
%         المصادر: القسم 4.2 (عقود الوكلاء)، backend/app/agents/،
%         نتائج القياس في evaluation/
%   5.2.2 الواجهة الخلفية والتهيئة الذكية وسجل التدقيق —
%         المالك: عبدالكريم العبدلي — 24/11 → 14/12
%         (مهمتا جانت ٢.٣.٥ التهيئة الذكية للربط والفهرسة،
%          و٢.٣.٦ سجل التدقيق والصلاحيات)
%         المصادر: الأقسام 4.3 و4.5 و4.6، وbackend/app/{api,db,core}/
%   5.2.3 تنفيذ الواجهة الأمامية — المالك: محمد الشريف — 01/12 → 10/12
%         (مهمة جانت د٣ ٢.٣.٣: دمج الوكلاء وبناء الواجهة)
%         المصادر: القسم 4.7 (الواجهات والمخططات السلكية)، frontend/
% قاعدة عامة: لا يوثَّق في هذا القسم إلا ما نُفّذ فعلاً في الكود، وتُرفق
%   لقطات شاشة ومخططات تنفيذ حقيقية لا مفاهيمية
%------------------------------------------------------------------------------
\section{تنفيذ النظام}
\label{sec:implementation}
~

%------------------------------------------------------------------------------
% 5.3 النشر والتشغيل
% المالك: وليد العنسي — 05/12 → 10/12 (يكتب مع مهمة النشر في ٢.٣.٨)
% المصادر المعتمدة:
%   (1) قرار البنية التحتية ADR-05 (محليًا أولًا + GPU بالساعة للتقييم)
%   (2) docker-compose.yml الإنتاجي ومجلد infra/ في مستودع الكود
%   (3) القسم 4.5 (ضوابط الأمان عند النشر) وNFR-08/NFR-09
% المخرجات المتوقعة: مخطط نشر TikZ، خطوات النشر والتهيئة، إدارة الأسرار
%   والمتغيرات البيئية، النسخ الاحتياطي والاسترجاع، مراقبة التشغيل
%------------------------------------------------------------------------------
\section{النشر والتشغيل}
\label{sec:deployment}
~

%------------------------------------------------------------------------------
% 5.4 اختبار النظام — المالك: محمد نجيب العواضي
% الموعد: كتابة تدريجية 15/11 → 10/12 (مهمة جانت ٢.٣.٧ المستمرة)،
%   ثم استكمال النتائج النهائية 09/12 → 14/12 (مهمة د٤ ٢.٣.٨)
% المصادر المعتمدة:
%   (1) إطار الاختبارات في مستودع الكود: backend/tests/ (وحدة/أدوات/
%       تكامل/e2e/أمان) وإطار التقييم evaluation/ (مجموعات الأسئلة
%       مع مخرجاتها المتوقعة وبرامج التشغيل والمقارنة)
%   (2) معايير القبول: الفصل 3 (المتطلبات وحالات الاستخدام) ومتطلبات
%       الجودة NFR المقابلة
%   (3) أهداف خطة التنفيذ: دقة ≥ 90% وزمن استجابة p95 ≤ 3 ثوانٍ
% المخرجات المتوقعة: منهجية الاختبار وأدواته، ثم نتائج فعلية مرقمة لكل
%   نوع (وحدة/تكامل/أداء/قبول) بجداول ومقارنة مع المستهدف
%------------------------------------------------------------------------------
\section{اختبار النظام}
\label{sec:testing}
~

\subsection{اختبارات الوحدة (\EN{Unit Tests})}
\label{subsec:unit-tests}
~

\subsection{اختبارات التكامل (\EN{Integration Tests})}
\label{subsec:integration-tests}
~

\subsection{اختبارات الأداء}
\label{subsec:performance-tests}
~

\subsection{اختبار قبول المستخدم (\EN{UAT})}
\label{subsec:uat}
~

%------------------------------------------------------------------------------
% 5.5 الخطة الزمنية للتنفيذ — المالك: أحمد المتوكل — 10/12 → 14/12
% المصادر: مهمات جانت ٢.١ → ٢.٣ (المرحلة الثانية)، سجل الدورات الفعلية
%   (طلبات السحب في المستودعين ونتائج القياس الأسبوعية)، المعالم M8 → M10
% المخرجات: جدول مقارنة المخطط بالمنفذ (الفروق وأسبابها) وملخص الدورات
%   التطويرية الأربع وما تحقق في كل منها
%------------------------------------------------------------------------------
\section{الخطة الزمنية للتنفيذ}
\label{sec:impl-schedule}
~

%------------------------------------------------------------------------------
% 5.6 ملخص الفصل — المالك: أحمد المتوكل — 12/12 → 14/12
% يُكتب أخيرًا بعد اكتمال 5.1 → 5.5 واعتماد نتائج الاختبار، ويُختم بمهام
%   المراجعة الختامية (مهمة جانت ٢.٣.٩: مراجعة الفصل الخامس واعتماده)
%------------------------------------------------------------------------------
\section{ملخص الفصل}
\label{sec:impl-summary}
~
